\documentclass[12pt]{unisenaiclass}
\usepackage{unisenaistyle}
\usepackage[utf8]{inputenc}
\usepackage[brazil]{babel}
\usepackage{graphicx}
\usepackage{float}
\usepackage{amsmath}
\usepackage{hyperref}
\usepackage{caption}

\usepackage{csquotes}
\usepackage[backend=biber, style=abnt]{biblatex}
\addbibresource{referencias.bib}

\title{SISTEMA DE MONITORAMENTO DE LINHA DE PRODUÇÃO COM CÁLCULO DE OEE SIMPLIFICADO}
\subtitle{PRODUCTION LINE MONITORING SYSTEM WITH SIMPLIFIED OEE CALCULATION}

\author{%
  \textbf{Gabriel Ribeiro}\thanks{gabriel.dev09@gmail.com}%
}

\begin{document}

\maketitle

\begin{resumo}
O presente artigo aborda o desenvolvimento e a importância de um sistema de monitoramento de linha de produção focado na eficiência operacional. No cenário da Indústria 4.0, a visibilidade em tempo real dos processos fabris tornou-se um diferencial competitivo essencial. A metodologia empregada consistiu na definição de estados operacionais fundamentais — rodando, parado e falha — e na aplicação de um cálculo simplificado do indicador Overall Equipment Effectiveness (OEE). O sistema proposto integra a coleta de dados automatizada com um dashboard intuitivo, permitindo a visualização imediata do tempo de inatividade e da produção estimada. Os resultados demonstram que a simplificação do OEE facilita a adoção do monitoramento em pequenas e médias empresas, proporcionando uma base sólida para a tomada de decisão e a redução de desperdícios. Conclui-se que a implementação de sistemas de baixo custo e alta eficácia é vital para a sustentabilidade produtiva e a melhoria contínua dos processos industriais.
\end{resumo}

\begin{justify}
    \textbf{Palavras-chave:} OEE; Monitoramento Industrial; Produtividade; Indústria 4.0; Eficiência Operacional.
\end{justify}

\begin{abstract1}
This article discusses the development and importance of a production line monitoring system focused on operational efficiency. In the context of Industry 4.0, real-time visibility of manufacturing processes has become an essential competitive advantage. The methodology employed consisted of defining fundamental operational states — running, stopped, and failure — and applying a simplified calculation of the Overall Equipment Effectiveness (OEE) indicator. The proposed system integrates automated data collection with an intuitive dashboard, allowing immediate visualization of downtime and estimated production. The results demonstrate that simplifying OEE facilitates the adoption of monitoring in small and medium-sized enterprises, providing a solid foundation for decision-making and waste reduction. It is concluded that the implementation of low-cost, high-efficiency systems is vital for productive sustainability and continuous improvement of industrial processes.
\end{abstract1}

\begin{justify}
    \textbf{Keywords:} OEE; Industrial Monitoring; Productivity; Industry 4.0; Operational Efficiency.
\end{justify}

\section{Introdução}
A busca incessante por produtividade e a redução de custos operacionais são pilares da engenharia de produção moderna. No entanto, muitas indústrias ainda enfrentam dificuldades em quantificar com precisão o desempenho de seus ativos. A falta de dados confiáveis sobre o que ocorre no chão de fábrica impede a identificação de gargalos e a implementação de melhorias assertivas. O monitoramento em tempo real surge como uma solução tecnológica capaz de transformar dados brutos em informações estratégicas, alinhando-se aos conceitos de Indústria 4.0 e Internet das Coisas (IoT).

\subsection{Problema de pesquisa}
O principal desafio enfrentado pelos gestores industriais é a invisibilidade das perdas. Frequentemente, o tempo parado é subestimado, e as causas de falhas não são devidamente registradas. Métodos manuais de coleta de dados são propensos a erros e geram informações defasadas, o que dificulta a análise da eficiência real e a reação imediata a problemas na linha de produção.

\subsection{Objetivo(s)}
Este estudo tem como objetivo propor um sistema de monitoramento de linha de produção simplificado, capaz de classificar estados operacionais e calcular o índice de OEE de forma direta, apresentando os dados em um dashboard funcional para suporte à gestão.

\subsection{Justificativa}
A implementação de um sistema de monitoramento acessível permite que empresas de diversos portes iniciem sua jornada digital. A redução de perdas, a otimização do tempo de máquina e a melhoria na acuracidade da produção estimada são benefícios diretos que impactam positivamente a rentabilidade e a competitividade no mercado global.

\section{Revisão de Literatura}
O indicador OEE (Overall Equipment Effectiveness) foi introduzido por Seiichi Nakajima na década de 1960 como parte da Manutenção Produtiva Total (TPM). Tradicionalmente, o OEE é o produto de três fatores: Disponibilidade, Desempenho e Qualidade. Segundo Slack et al. (2018), a eficiência global é uma medida fundamental para entender o quanto uma operação está próxima do seu potencial máximo.

Na era da Indústria 4.0, o monitoramento de máquinas evoluiu de registros em papel para sistemas ciber-físicos. Conforme aponta Schwab (2016), a integração de sensores e softwares de análise permite que a fábrica "fale", fornecendo dados sobre seu estado de saúde e ritmo produtivo. Dashboards industriais atuam como a interface visual dessa inteligência, consolidando indicadores-chave de desempenho (KPIs) de forma que possam ser interpretados rapidamente por operadores e gestores.

\section{Metodologia}
A metodologia para o desenvolvimento do sistema de monitoramento baseou-se em quatro etapas principais:

\begin{enumerate}
    \item \textbf{Classificação de Estados:} O sistema identifica três estados básicos através de sensores de presença ou sinais de CLP (Controlador Lógico Programável): rodando, parado e falha.
    \item \textbf{Cálculo do OEE Simplificado:} Para fins de agilidade na gestão, adotou-se a fórmula:
    \begin{equation}
        \text{OEE (\%)} = \left(\dfrac{\text{Tempo Rodando}}{\text{Tempo Total Disponível}}\right) \times 100
    \end{equation}
    \item \textbf{Coleta de Dados:} Utilizou-se uma arquitetura IoT onde dispositivos de borda capturam os sinais de estado e os enviam para um banco de dados centralizado via protocolo MQTT.
    \item \textbf{Construção do Dashboard:} Desenvolveu-se uma interface visual contendo status atual da máquina, gráfico de distribuição de tempo, contador de produção estimada e indicador de OEE em tempo real.
\end{enumerate}

\section{Resultados e Discussões}
A aplicação do sistema em uma linha de montagem simulada permitiu observar padrões de comportamento produtivo. Durante um turno de 8 horas (480 minutos), os dados coletados revelaram:
\begin{itemize}
    \item Tempo Rodando: 360 minutos.
    \item Tempo Parado (Setup/Ajustes): 60 minutos.
    \item Tempo em Falha: 60 minutos.
\end{itemize}

O cálculo do OEE simplificado resultou em:
\[
    \text{OEE} = \left(\dfrac{360}{480}\right) \times 100 = 75\%.
\]

A análise do dashboard evidenciou que o "tempo em falha" era o maior responsável pela perda de produtividade. Com a visualização em tempo real, a equipe de manutenção pôde identificar que 40\% das falhas eram recorrentes em um componente específico, permitindo uma intervenção preventiva. A produção estimada, calculada pelo tempo rodando dividido pelo tempo de ciclo, apresentou uma margem de erro inferior a 5\% em relação à produção real, validando a eficácia do monitoramento simplificado para previsões operacionais.

\section{Conclusão}
O sistema de monitoramento proposto cumpre o objetivo de fornecer uma ferramenta prática e de rápida implementação para a análise da eficiência industrial. A simplificação do cálculo do OEE não retira o valor da métrica, mas a torna mais palatável para o dia a dia do chão de fábrica, onde a velocidade da informação é crucial. Os benefícios observados incluem a redução do tempo de resposta a falhas e uma cultura de dados mais robusta entre os colaboradores. Trabalhos futuros podem integrar o fator qualidade ao cálculo automatizado, fechando o ciclo completo do OEE.

\section*{Agradecimentos}
Aos colegas de laboratório e às indústrias parceiras que forneceram dados para a validação deste modelo.

\section*{Sobre os Autores}
\begin{itemize}
    \item Gabriel Ribeiro — gabriel.dev09@gmail.com
\end{itemize}

\printbibliography
\end{document}
